\chapter{Hitchin 系统的量子化}\label{ch:14}
% 原PDF第67页，第14章开始。
本章旨在解释 Hitchin Hamilton 量的 Beilinson–Drinfeld 量子化，以及由此得到的 Hecke 特征层\cite{BD91,BD05}。仍记自守群为 $G$，其 Langlands 对偶群为 $\checkG$，并假设 $\checkG$ 为伴随型。记 $\check{\mathfrak g}=\operatorname{Lie}(\checkG)$，固定 Borel 子群 $\check B\subset\checkG$，置
\[\mathfrak c_{\checkG}:=\check{\mathfrak g}/\!/\checkG=\Spec k[\check{\mathfrak g}]^{\checkG},\qquad
\mathfrak c_{\checkG,\Omega}:=(\mathfrak c_{\checkG})_{\Omega_X}.\]
若 $d_1,\ldots,d_r$ 为 $\check{\mathfrak g}$ 的指数，则 $\mathfrak c_{\checkG,\Omega}\simeq\bigoplus_i\Omega_X^{d_i+1}$，Hitchin 底空间为
\[\operatorname{Hitch}_{\checkG}(X):=\Gamma(X,\mathfrak c_{\checkG,\Omega}).\]
本章的要点是，同一个仿射空间出现两次：在经典情形，它是 Hitchin 底空间；在量子情形，它是临界水平中心的相伴分次。

\section{oper 作为 Hitchin 底空间上的挠子}\label{sec:14.1}\index{oper|hyperpage}\index{Hitchin 底空间|hyperpage}\index{射流|hyperpage}
\srcpage{68}
令 $\pi:Y\to X$ 为 $D_X$-概形。
\begin{definition*}{}
依照 Beilinson–Drinfeld\cite{BD05}，$Y$ 上相对于 $X$ 的 $\checkG$-oper 是一个二元组 $(\mathcal F_{\check B},\nabla)$，其中 $\mathcal F_{\check B}$ 是 $Y$ 上的 $\check B$-挠子，$\nabla$ 是诱导的 $\checkG$-挠子 $\mathcal F_{\checkG}:=\checkG\times^{\check B}\mathcal F_{\check B}$ 上相对于 $\pi$ 的联络，并且 $\check B$-约化的第二基本形式满足 oper 条件：
\[c(\nabla)\in\Gamma(Y,(\check{\mathfrak g}/\check{\mathfrak b})_{\mathcal F_{\check B}}\otimes\pi^*\Omega_X)\]
位于由简单负根方向组成的子丛中的开 $\check B$-轨道；等价地，它的简单负根分量处处非零。
\end{definition*}
函子 $Y\mapsto\Op_{\checkG}(Y)$ 由一个 $D_X$-概形表示，记为 $\Op^D_{\checkG}\to X$。它的全局截面是通常的 oper 空间 $\Op_{\checkG}(X)$。
\begin{proposition}{}{14.1.1}
预叠 $\Op_{\checkG}(X)$ 是向量空间 $\operatorname{Hitch}_{\checkG}(X)$ 的挠子。因此，若

\begin{transnote}{含糊点拨}
底空间 $\operatorname{Hitch}_{\checkG}(X)=\bigoplus_i H^0(X,\Omega_X^{d_i+1})$ 是仿射空间，$d_i$ 是 $\check{\mathfrak g}$ 齐次生成元的次数；oper 空间是它的挠子：任取两个 oper 之差给出一个 Hitchin 场，但没有典范的 oper 可以固定，所以只能说是「挠子」。过滤来自 Rees 构造，因此 $\operatorname{gr} A\simeq A^{\mathrm{cl}}$——这就是「量子化在经典极限回到 Hitchin 系统」的确切含义。[来源：nLab "Hitchin fibration"（该条目无 Last revised，2026-11-17 访问）；并见原书 p.68 本段。]
\end{transnote}
\[A:=\Gamma(\Op_{\checkG}(X),\cO),\qquad A^{\mathrm{cl}}:=\Gamma(\operatorname{Hitch}_{\checkG}(X),\cO),\]
则 $A$ 的自然过滤的相伴分次为 $\operatorname{gr}A\simeq A^{\mathrm{cl}}$。
\end{proposition}

命题\ref{pro:14.1.1}的证明归约到引理\ref{lem:14.1.2}、\ref{lem:14.1.3}和\ref{lem:14.1.4}。
\begin{lemma}{}{14.1.2}
$\Op_{\checkG}(X)$ 非空。
\end{lemma}
\Needspace{8\baselineskip}
\begin{lemma}{}{14.1.3}
令 $W$ 为遗忘函子 $\operatorname{Sch}_{/X_{\mathrm{dR}}}\to\operatorname{Sch}_{/X}$ 的右伴随。对 $X$ 上的向量丛 $E$，可以将 $W(E)$ 理解为射流的 $D_X$-概形，使得
\[W(E)(D_x)=\operatorname{Jets}_x(E).\]
则 $\Op^D_{\checkG}$ 是 $D_X$-概形 $W(\mathfrak c_{\checkG,\Omega})$ 的挠子。
\end{lemma}
\begin{lemma}{}{14.1.4}
在 $\check{\mathfrak g}$ 中选定主 $\mathfrak{sl}_2$ 三元组 $(e,h,f)$。存在依赖于这一选择的 $X$ 上典范 $\check B$-挠子 $\mathcal F^{\mathrm{op}}_{\check B}$，使得对于任意 $Y\to X$ 上的每个 $\checkG$-oper，其底层 $\check B$-挠子都典范地等于 $\mathcal F^{\mathrm{op}}_{\check B}$ 的拉回。相应的 $\Bun_{\checkG}$ 中的点，在连通分支的意义下与辅助选择无关。
\end{lemma}

特别地，遗忘联络给出到 $\Bun_{\checkG}$ 的常值映射：
\origdiagram{p068-square.pdf}{0.3484}

\section{局部正规形}\label{sec:14.2}\index{Drinfeld–Sokolov 规范变换|hyperpage}\index{Kostant 切片|hyperpage}
在 $X$ 上局部讨论，选择坐标 $t$，并先假设典范的 $\check B$-挠子平凡。在这样的平凡化下，oper 联络具有形式
\[\nabla=d+(f+\rho(t))\,dt,\qquad\rho(t)\in\check{\mathfrak b}((t)),\]
正则情形下则有 $\rho(t)\in\check{\mathfrak b}[[t]]$。这种形式并不唯一。oper 的 Drinfeld–Sokolov 规范变换陈述\cite{BD05,Fre02}说，每个 $\check N((t))$-规范等价类都有唯一的代表元
\[\nabla=d+(f+v(t))\,dt,\qquad v(t)\in\check{\mathfrak g}^{e}((t)),\]
正则 oper 满足 $v(t)\in\check{\mathfrak g}^{e}[[t]]$。这里 $\check{\mathfrak g}^{e}$ 是 $e$ 的中心化子。这就是 Kostant 切片：真正的参数是 $\check{\mathfrak g}^{e}$ 中的 Kostant 坐标，而非整个 $\check{\mathfrak b}$。

\begin{transnote}{含糊点拨}
$\check N((t))$ 以 $n(t)\mapsto\nabla+n(t)$ 的方式作用在 oper 上；Kostant 的 $\mathfrak{sl}_2$ 理论给出横截切片：$\check{\mathfrak g}^{e}$ 与 $\check{\mathfrak b}$ 交于唯一代表。$f$ 定 oper 的最低次项，$e$ 定切片方向，因此真正的参数是 $\check{\mathfrak g}^{e}$ 里的 Kostant 坐标，而不是整个 $\check{\mathfrak b}$——这就是「Kostant 切片」在此的含义。[来源：nLab "Drinfeld-Sokolov reduction"，https://ncatlab.org/nlab/show/Drinfeld-Sokolov+reduction（Last revised 2013-06-10；2026-11-17 访问）；并见原书 p.68 本段与 [BD05]、[Fre02]。]
\end{transnote}
\srcpage{69}
\begin{proof}

先取 $\checkG=\operatorname{PGL}_2$。oper 条件平凡化 $(\check{\mathfrak g}/\check{\mathfrak b})_{\mathcal F_{\check B}}\otimes\pi^*\Omega_X$。因此，$\Op_{\operatorname{PGL}_2}(Y)$ 是 $\pi^*\Omega_X^{\otimes2}$ 的一个可能为空的挠子。当 $g>1$ 时，由于 $H^1(X,\Omega_X^{\otimes2})=0$，它非空；当 $g=0,1$ 时，可以直接构造。因此，三个引理对 $\operatorname{PGL}_2$ 成立。在这一情形中，通常的陈述是：射影联络组成二次微分空间的挠子。

对于一般伴随型 $\checkG$，沿主同态 $\operatorname{PGL}_2\to\checkG$ 诱导一个 $\operatorname{PGL}_2$-oper。这给出一个 $\checkG$-oper，从而证明非空性。Kostant 同构给出
\[(\check{\mathfrak g}^{e})_{\mathcal F^{\mathrm{op}}_{\check B}}\otimes\Omega_X\simeq\mathfrak c_{\checkG,\Omega},\]
所以同一论证将一般 oper 的自由参数识别为 Hitchin 向量丛。这证明 $\Op^D_{\checkG}$ 是 $W(\mathfrak c_{\checkG,\Omega})$-挠子，进而 $\Op_{\checkG}(X)$ 是 $\operatorname{Hitch}_{\checkG}(X)$-挠子。
\end{proof}

\section{Feigin–Frenkel 中心}\label{sec:14.3}\index{Feigin–Frenkel 中心|hyperpage}\index{临界水平|hyperpage}\index{真空模|hyperpage}\index{手征代数|hyperpage}
令 $x\in X$，$D_x=\Spec k[[t]]$，$D_x^\times=\Spec k((t))$。记 $\widehat{\mathfrak g}_{\mathrm{crit},x}$ 为临界水平的仿射 Kac–Moody 代数，
\[\mathbb V_{\mathrm{crit},x}:=\operatorname{Ind}_{\mathfrak g[[t]]\oplus k\mathbf1}^{\widehat{\mathfrak g}_{\mathrm{crit},x}}k\]
为真空模。有两个相关的中心：
\[Z_{\mathrm{crit},x}:=Z(\widehat U_{\mathrm{crit},x}),\qquad
\mathfrak z_{\mathrm{crit}}(D_x):=\End_{\widehat{\mathfrak g}_{\mathrm{crit},x}}(\mathbb V_{\mathrm{crit},x}).\]
真空模随 $x$ 变动，组成临界手征代数 $\mathbb V_{\mathrm{crit}}$。它的中心 $\mathfrak z_{\mathrm{crit}}:=Z(\mathbb V_{\mathrm{crit}})$ 是交换手征代数；等价地，它是因子化意义下的交换 $D_X$-代数，所以 $\Spec\mathfrak z_{\mathrm{crit}}$ 是一个 $D_X$-概形。
\begin{theorem}{Feigin–Frenkel，参见\cite{Fre02}}{14.3.1}
存在典范识别
\[\operatorname{Spf}Z_{\mathrm{crit},x}\simeq\Op_{\checkG}(D_x^\times),\qquad
\Spec\mathfrak z_{\mathrm{crit}}(D_x)\simeq\Op_{\checkG}(D_x),\]
在全局，有 $D_X$-概形的识别
\[\Spec\mathfrak z_{\mathrm{crit}}\simeq\Op^D_{\checkG}.\]
在 PBW 过滤下，经典极限是
\[\Spec\operatorname{gr}\mathfrak z_{\mathrm{crit}}\simeq W(\mathfrak c_{\checkG,\Omega}).\]
\end{theorem}
\begin{proof}

证明思路如下。首先，$\mathbb V_{\mathrm{crit}}$ 的手征代数结构使 $\mathfrak z_{\mathrm{crit}}$ 具有因子化性质；在局部，这产生一个从通常临界中心的形式谱到穿孔圆盘中心的映射。\thref{14.3.1}说明这个映射是同构。因此，剩下的任务是将 $\Spec\mathfrak z_{\mathrm{crit}}(D)$ 与 $D$ 上的正则 oper 识别，并使这一识别与 $\Aut(D)$ 相容。

正则映射 $\Spec\mathfrak z_{\mathrm{crit}}(D)\to\Op_{\checkG}(D)$ 通过 Tannaka 方法构造。需要在 $\Spec\mathfrak z_{\mathrm{crit}}(D)\times D$ 上构造一个 $\checkG$-挠子、一个联络，以及满足 oper 条件的 $\check B$-约化，并要求这些数据都是 $\Aut(D)$-等变的。几何 Satake 与临界全局截面给出对称幺半函子
\[F_D:\Rep(\checkG)\longrightarrow\mathfrak z_{\mathrm{crit}}(D)\text{-}\mathrm{Mod},\]
由性质
\[\Gamma(\Gr_G,\operatorname{Sat}(V))_{\mathrm{crit}}\simeq F_D(V)\otimes_{\mathfrak z_{\mathrm{crit}}(D)}\mathbb V_{\mathrm{crit}}\]
刻画。由 Tannaka 对偶，$F_D$ 确定 $\Spec\mathfrak z_{\mathrm{crit}}(D)$ 上的 $\checkG$-挠子。$\check B$-约化从最低旋转权得到。对不可约 $\checkG$-表示 $V^{\check\lambda}$，令 $\operatorname{IC}_{\check\lambda}=\operatorname{Sat}(V^{\check\lambda})$。若 $L_0=-t\partial_t\in\operatorname{Lie}\Aut^0(D)$，则 $\Gamma(\Gr_G,\operatorname{IC}_{\check\lambda})_{\mathrm{crit}}$ 的最低 $L_0$-特征值为 $-\langle\check\lambda,\rho\rangle$；相应特征空间是一维的，并且 $\check B$-不变。这些直线满足 Plücker 关系，从而给出所需的 $\check B$-约化。最后一步是证明得到的到 $\Op_{\checkG}(D)$ 的映射为同构。

\begin{transnote}{含糊点拨}
$L_0=-t\partial_t$ 来自形式圆盘的旋转（伸缩）作用；最低特征值 $-\langle\check\lambda,\rho\rangle$ 里的 $\rho$ 是半和正根，这正是临界扭转与 $\rho$-移位的来源。最低权空间一维且 $\check B$-不变，才能当纤维化的一根直线；这些直线之间满足 Plücker 关系，来自 $\mathrm{IC}_{\check\lambda}=\operatorname{Sat}(V^{\check\lambda})$ 的函子性，而不是额外构造。[依据：原书 p.69 本段；并见 [Fre02]。]
\end{transnote}
\end{proof}
\section{全局微分算子与量子 Hitchin 映射}\label{sec:14.4}\index{量子 Hitchin 映射|hyperpage}\index{临界扭转微分算子|hyperpage}
\srcpage{70}
令 $\cD_{\mathrm{crit}}$ 为 $\Bun_G$ 上临界扭转微分算子层。局部临界中心通过局部化作用于全局微分算子。在局部，使用如下标准构造：若一对 $(\mathfrak h,K)$ 作用于光滑叠 $Y$，而 $y\in Y$ 是 $K$-不动点，则齐性模型上的不变微分算子映入 $Y$ 上的微分算子：
\[(U(\mathfrak h)/U(\mathfrak h)\mathfrak k)^K\longrightarrow\Gamma(Y,\cD_Y).\]
取 $Y=\Bun_G$、$\mathfrak h=\mathfrak g((t))$、$K=G[[t]]$，得到
\[\mathfrak z_{\mathrm{crit}}(D_x)=\End(\mathbb V_{\mathrm{crit},x})\longrightarrow\Gamma(\Bun_G,\cD_{\mathrm{crit}}).\]
这些局部映射与变量 $x$ 中的联络相容，因此下降到水平截面：
\[\mathfrak z_{\mathrm{crit}}(X)\longrightarrow\Gamma(\Bun_G,\cD_{\mathrm{crit}}).\]
取相伴分次给出经典 Hitchin 映射
\[\operatorname{gr}\mathfrak z_{\mathrm{crit}}(X)\longrightarrow\Gamma(T^*\Bun_G,\cO).\]

利用 Feigin–Frenkel 识别与上述挠子陈述，置
\[\begin{aligned}
A&:=\Gamma(\Op_{\checkG}(X),\cO)=\Gamma(X_{\mathrm{dR}},\mathfrak z_{\mathrm{crit}}),\\
A^{\mathrm{cl}}&:=\Gamma(\operatorname{Hitch}_{\checkG}(X),\cO)\simeq\operatorname{gr}A.
\end{aligned}\]
\begin{theorem*}{}
上述映射为同构\cite{BD91}：
\[A\xrightarrow{\sim}\Gamma(\Bun_G,\cD_{\mathrm{crit}}),\qquad
A^{\mathrm{cl}}\xrightarrow{\sim}\Gamma(T^*\Bun_G,\cO).\]
\end{theorem*}
对 $\GL_n$ 方向上的正特征方法，参见\cite{BCTZ16}。因此，量子 Hitchin Hamilton 量的交换代数，就是 $\checkG$-oper 空间上的函数代数。第二个同构是经典陈述：$T^*\Bun_G$ 上的全局函数来自 Hitchin 底空间。

\section{oper 对应的自守模}\label{sec:14.5}\index{Hecke 特征层!由 oper 构造|hyperpage}\index{Beilinson–Drinfeld 局部化|hyperpage}
令 $\sigma\in\Op_{\checkG}(X)$。通过 $\Op_{\checkG}(X)\subset\LocSys_{\checkG}(X)$，它有一个底层 $\checkG$-局部系统，仍记为 $\sigma$。令 $k_\sigma$ 为 $A=\Gamma(\Op_{\checkG}(X),\cO)$ 的相应特征标。定义临界扭转 $D$ 模
\begin{definition*}{}
\[\Aut_\sigma:=\cD_{\mathrm{crit}}\otimes_{A,\sigma}k_\sigma\in\DMod(\Bun_G).\]
\end{definition*}
\begin{theorem}{Beilinson–Drinfeld\cite{BD91}}{14.5.1}
$\Aut_\sigma$ 是特征值为 $\sigma$ 的 Hecke 特征层。等价地，对每个 $x\in X$ 及每个 $V\in\Rep(\checkG)$，
\[H_{V,x}(\Aut_\sigma)\simeq V_{\sigma,x}\otimes\Aut_\sigma,\]
其中 $H_{V,x}$ 是对应于 $V$ 的 Hecke 函子，而 $V_{\sigma,x}$ 是由 $V$ 和 $\sigma$ 关联的局部系统在 $x$ 处的纤维。
\end{theorem}
\begin{proof}

证明使用 Beilinson–Drinfeld 局部化
\[\operatorname{Loc}_x:\widehat{\mathfrak g}_{\mathrm{crit},x}\text{-}\mathrm{Mod}^{G[[t]]}\longrightarrow\DMod_{\mathrm{crit}}(\Bun_G),\]
它满足
\[\operatorname{Loc}_x(\mathbb V_{\mathrm{crit},x})\simeq\cD_{\mathrm{crit}}.\]
\srcpage{71}
对 $V\in\Rep(\checkG)$，Hecke 作用与局部化相容：
\[\begin{aligned}
\operatorname{Sat}_x(V)\star\cD_{\mathrm{crit}}
&\simeq\operatorname{Loc}_x(\operatorname{Sat}_x(V)\star\mathbb V_{\mathrm{crit},x})\\
&\simeq\operatorname{Loc}_x\bigl(F_{D_x}(V)\otimes_{\mathfrak z_{\mathrm{crit}}(D_x)}\mathbb V_{\mathrm{crit},x}\bigr).
\end{aligned}\]
令 $\mathcal F$ 为 $\Op_{\checkG}(X)$ 上的普遍 $\checkG$-挠子，等价地，是普遍 oper 局部系统。函子 $F_{D_x}$ 是这一普遍挠子在形式圆盘上的限制，所以上式变为
\[\operatorname{Sat}_x(V)\star\cD_{\mathrm{crit}}\simeq V_{\mathcal F,x}\otimes_A\cD_{\mathrm{crit}}.\]
在 $A$ 上与 $k_\sigma$ 作张量积，得到
\[\operatorname{Sat}_x(V)\star\Aut_\sigma\simeq V_{\sigma,x}\otimes\Aut_\sigma.\]
这正是 Hecke 特征性质。
\end{proof}

\section{真空生成与下一步的局部图景}\label{sec:14.6}\index{基本局部等价|hyperpage}\index{非分歧 oper|hyperpage}
\secref{14.5}的构造只使用临界 Kac–Moody 范畴中由真空模生成的部分。示意地，有
\origdiagram{p071-square.pdf}{0.4918}
其中，$\mathbb V$-生成意为由真空对象生成。下一步是用整个临界范畴替换这个真空生成的子范畴。预期的局部形式是基本局部等价。置
\[\Op^{\mathrm{unr}}_{\checkG,x}:=\Op_{\checkG}(D_x^\times)\times_{\LocSys_{\checkG}(D_x^\times)}\LocSys_{\checkG}(D_x).\]
于是，忽略通常的 $*$ 型与 $!$ 型 $\IndCoh$ 约定的区别，
\[\operatorname{KL}_{G,\mathrm{crit},x}\simeq\IndCoh(\Op^{\mathrm{unr}}_{\checkG,x}).\]
局部化应当提升为如下形式的相对构造：
\[\begin{split}
\operatorname{KL}_{G,\mathrm{crit},x}\otimes_{\QCoh(\Op^{\mathrm{unr}}_{\checkG,x})}
\QCoh(\Op_{\checkG}(X\setminus\{x\})^{\mathrm{unr}})\\
\longrightarrow\DMod_{\mathrm{crit}}(\Bun_G),
\end{split}\]
其中定义
\[\Op_{\checkG}(X\setminus\{x\})^{\mathrm{unr}}:=\Op_{\checkG}(X\setminus\{x\})\times_{\LocSys_{\checkG}(X\setminus\{x\})}\LocSys_{\checkG}(X).\]
这就是从 oper 的 Beilinson–Drinfeld 特征层，通向\cite{FG05,FG07}中范围更广的临界水平局部理论的桥梁。

